\subsection{Conjunction}

\begin{notation}
    \[P \wedge Q \coloneqq \neg P \nor \neg Q\]
\end{notation}

% P, Q -> P and Q
\begin{align*}
    \premise{P}\\
    \premise{Q}\\
    \assumption[1]{\neg P}\\
    \byrule[1]{\bot}{\negE} \\
    \assumption[1]{\neg Q}\\
    \byrule[1]{\bot}{\negE}\\
    \byrule{\neg P \nor \neg Q}{\norI}\\
    \bydefinition{P \wedge Q}
\end{align*}
\newcommand{\wedgeI}{{\wedge}I}
\[\inferrule{P \\ Q}{P \wedge Q} \quad (\wedgeI)\]

% P and Q -> P, Q
\begin{align*}
    \premise{P \wedge Q}\\
    \bydefinition{\neg P \nor \neg Q}\\
    \assumption[1]{\neg P}\\
    \byrule[1]{\bot}{\norE}\\
    \byrule{P}{\raa}
\end{align*}
\begin{align*}
    \premise{P \wedge Q}\\
    \bydefinition{\neg P \nor \neg Q}\\
    \assumption[1]{\neg Q}\\
    \byrule[1]{\bot}{\norE}\\
    \byrule{Q}{\raa}
\end{align*}
% Therefore, wizardry?
\newcommand{\wedgeE}{{\wedge}E}
\[\inferrule{P \wedge Q}{P} \quad \inferrule{P \wedge Q}{Q} \quad (\wedgeE)\]

%\begin{align*}
%    \premise{P \wedge Q}\\
%    \byrule{Q}{\wedgeE}\\
%    \byrule{P}{\wedgeE}\\
%    \byrule{Q \wedge P}{\wedgeI}
%\end{align*}
%%\newcommand{\wedgeC}{{\wedge}C}
%\[\inferrule{P \wedge Q}{Q \wedge P} \quad (\wedgeC)\]

%\begin{align*}
%    \premise{(P \wedge Q) \wedge R}\\
%    \byrule{P \wedge Q}{\wedgeE}\\
%    \byrule{P}{\wedgeE}\\
%    \byrule{Q}{\wedgeE}\\
%    \byrule{R}{\wedgeE}\\
%    \byrule{Q \wedge R}{\wedgeI}\\
%    \byrule{P \wedge (Q \wedge R)}{\wedgeI}
%\end{align*}
%%\newcommand{\wedgeA}{{\wedge}A}
%\[\inferrule{(P \wedge Q) \wedge R}{P \wedge (Q \wedge R)} \quad (\wedgeA)\]
